\documentclass[../main]{subfiles}
\begin{document}

\chapter{Presentación Máquinas Térmicas.}

\img{images/17/maquina-de-vapor.jpg}

\section*{Temas a Elegir (Máquinas Térmicas)}
\begin{itemize}
    \item Caldera
    \item Cámara Frigorífica
    \item Bomba de Calor (Aire Acondicionado)
    \item Máquina de Vapor (Ciclo de Rankine)
    \item Central Térmica de Ciclo Combinado
    \item Motor Stirling
    \item Motor de 2 Tiempos
\end{itemize}

\section*{Estructura de la Presentación (Contenido Requerido)}

La presentación debe abordar de manera concisa y precisa los siguientes 9 puntos para el tema seleccionado:

\begin{enumerate}
    \item \textbf{Principio de Funcionamiento:} Explicación fundamental del proceso termodinámico.
    \item \textbf{Ciclo Térmico:} Explicar detalladamente el ciclo termodinámico asociado (ej. Ciclo de Rankine, Ciclo de Carnot, Ciclo de Otto, etc.) \textemdash se pueden incluir diagramas $P$-$v$ o $T$-$s$.
    \item \textbf{Aspecto Visual (Imágenes):} Incluir un mínimo de 10 fotografías o esquemas técnicos relevantes (internos, externos, diagramas de flujo, etc.).
    \item \textbf{Historia:} Reseña de su fabricación, primeras aplicaciones e hitos de desarrollo.
    \item \textbf{Características Técnicas Clave:}
    \begin{itemize}
        \item \textbf{Consumo y Energía:} Tipos de combustible/energía y rangos de consumo típicos.
        \item \textbf{Potencia y Eficiencia:} Especificaciones de potencia nominal y eficiencia térmica ($\eta$).
        \item \textbf{Clasificación:} Según el principio de funcionamiento o uso (ej. Motor de combustión interna, máquina de flujo, ciclo de refrigeración, etc.).
    \end{itemize}
    \item \textbf{Mantenimiento:} Descripción de las tareas de mantenimiento preventivo y correctivo más críticas.
    \item \textbf{Ventajas y Desventajas:} Análisis comparativo respecto a tecnologías alternativas o modelos previos.
    \item \textbf{Mercado Actual:} Usos predominantes, tendencias de desarrollo e impacto económico-ambiental.
    \item \textbf{Conclusiones:} Síntesis del potencial, limitaciones e importancia del dispositivo en el panorama energético/industrial actual.
\end{enumerate}

\section*{Requerimiento Adicional (Fórmulas)}
Para el punto 2 (\textbf{Ciclo Térmico}) o 5 (\textbf{Características Técnicas}), se deben incluir las fórmulas termodinámicas fundamentales relacionadas con el rendimiento o el trabajo del ciclo.

\subsection*{Ejemplo de Fórmula (Eficiencia Térmica)}
La eficiencia térmica ($\eta$) de una máquina térmica genérica se define como la razón entre el trabajo neto producido ($W_{neto}$) y el calor suministrado ($Q_{sumin}$):
$$ \eta = \frac{W_{neto}}{Q_{sumin}} = 1 - \frac{Q_{cedido}}{Q_{sumin}} $$
Para una máquina de Carnot:
$$ \eta_{Carnot} = 1 - \frac{T_L}{T_H} $$
donde $T_L$ es la temperatura del foco frío y $T_H$ es la temperatura del foco caliente (en Kelvin).

\end{document}